\subsection{Post-trained Model Evaluations}
\label{subsec:final_model_evals}

\subsubsection{Evaluation Benchmarks}

\begin{table*}[!htp]
\centering
\small
\setlength{\tabcolsep}{7pt}
\renewcommand{\arraystretch}{1.15}

\begin{tabular}{l|c c c }
\toprule
\textbf{Benchmark} & \textbf{N-3-Nano} & \textbf{Qwen3} & \textbf{GPT-OSS} \\
\midrule

\rowcolor{black!5}
\multicolumn{4}{l}{\textbf{General Knowledge}} \\
MMLU-Pro &  78.30 &  \textbf{80.90} & 75.00 \\

\midrule
\rowcolor{black!5}
\multicolumn{4}{l}{\textbf{Reasoning}} \\
AIME25 (no tools) & 89.06 & 85.00 & \textbf{91.70} \\
AIME25 (with tools) & \textbf{99.17} & - & 98.7 \\
GPQA (no tools) &  73.04 & \textbf{73.40} & 71.50 \\
GPQA (with tools) & \textbf{75.00} & - & 74.20 \\
% LiveCodeBench (v5 2024-07$\leftrightarrow$12) & 72.063 & 66.000 & 78 \\
LiveCodeBench (v6 2024-08$\leftrightarrow$2025-05) & \textbf{68.25} & 66.00 & 61.00\\
SciCode (subtask) & 33.28 & 33.00 & \textbf{34.00}  \\
HLE (no tools) & 10.57 & 9.80 & \textbf{10.90} \\
HLE (with tools) & 15.48 & - & \textbf{17.30} \\
MiniF2F pass@1 & \textbf{50.03} & 5.72* & 12.05* \\
MiniF2F pass@32 & \textbf{79.92} & 16.80* & 43.03* \\
\midrule
\rowcolor{black!5}
\multicolumn{4}{l}{\textbf{Agentic}} \\
Terminal Bench (hard subset) & 8.51 & 5.00 & \textbf{10.00} \\
% Terminal Bench (nv-internal mix) &  &  &  \\
% SWE-Bench Verified & 38.760 & ? & 38.760 \\
SWE-Bench (OpenHands) & \textbf{38.76} & 22.00* & 34.00* \\
\textbf{TauBench V2} &  &  &  \\
\quad Airline & 48.00 & \textbf{58.00} & 38.00 \\
\quad Retail & 56.91 & \textbf{58.80} & 54.80 \\
\quad Telecom & 42.21 & 26.30 & \textbf{49.70} \\
\quad Average & \textbf{49.04} & 47.70 & 47.50 \\
% Aider-Polyglot &  &  &  \\
BFCL v4 & \textbf{53.76} & 46.40* & - \\
\midrule
\rowcolor{black!5}
\multicolumn{4}{l}{\textbf{Chat \& Instruction Following}} \\
IFBench (prompt) & \textbf{71.51} & 51.00 & 65.00 \\
% IFBench (instruction) & 77.015 &  &  \\
Scale AI Multi Challenge & 38.45  & \textbf{44.75} & 33.75 \\
Arena-Hard-V2 (Hard Prompt) & \textbf{72.10} & 49.60* & 71.20* \\
Arena-Hard-V2 (Creative Writing) & 63.20 &  \textbf{66.00*} & 25.90* \\
Arena-Hard-V2 (Average) & \textbf{67.65} & 57.80 & 48.55 \\

\midrule
\rowcolor{black!5}
\multicolumn{4}{l}{\textbf{Long Context}} \\
AA-LCR & 35.85 & \textbf{59.00} & 34.00 \\
RULER-100 @ 256k & \textbf{92.92} & 89.40 & - \\
RULER-100 @ 512K & \textbf{91.25} & 84.00 & - \\
RULER-100 @ 1M & \textbf{86.34} & 77.50 & - \\

\midrule
\rowcolor{black!5}
\multicolumn{4}{l}{\textbf{Multilingual}} \\
MMLU-ProX (avg over langs) & 59.50 & \textbf{77.60*} & 69.10* \\
% MMLU-ProX (avg over langs w/o Ja) &  &  &  \\
WMT24++ (en$\rightarrow$xx) & \textbf{86.20} & 85.60 & 83.20 \\
% COMET?? & &  &  \\

\bottomrule
\end{tabular}

\caption{
\ourmodel compared to \text{Qwen3-30B-A3B-Thinking-2507}, and \text{GPT-OSS 20B}. 
% \todo{JK: Structured Output environment} 
% \today{Alex: Genselect \& Formal proofs}
% \todo{Change precision to 1 decimal. There seems to be a discrepancy rn}
}
\label{tab:nano_v3_comparison}
\end{table*}


We evaluate \ourmodel across a broad suite of established benchmarks spanning mathematical and
scientific reasoning, coding, agentic tool use, instruction following, long-context understanding, and
multilingual capability. Table~\ref{tab:nano_v3_comparison} summarizes the final results.

All evaluation results were collected via Nemo Evaluator SDK\footnote{\url{https://github.com/NVIDIA-NeMo/Evaluator}} and for most benchmarks, the Nemo Skills Harness\footnote{\url{https://github.com/NVIDIA-NeMo/Skills}}. For reproducibility purposes, the open source container on Nemo Skills packaged via NVIDIA's Nemo Evaluator SDK used for evaluations can be found here\footnote{\url{https://catalog.ngc.nvidia.com/orgs/nvidia/teams/eval-factory/containers/nemo_skills}}. In addition to Nemo Skills, the evaluations also used dedicated packaged containers for Tau-2 Bench, ArenaHard v2, AA\_LCR. More details on the evaluation settings can be found in the Nemo Evaluator SDK configs folder\footnote{\url{https://github.com/NVIDIA-NeMo/Evaluator}}. The following benchmarks are not onboarded yet in our open source tools and for these we used their official open source implementation: Terminal Bench, SWE-Bench, Scale AI Multi Challenge.


For mathematical and STEM reasoning, we evaluate on \textsc{AIME25} (with and without tools), \textsc{GPQA}~\citep{rein2023gpqa}, \textsc{LiveCodeBench v6}~\citep{jain2024livecodebench},
\textsc{SciCode}~\citep{tian2024scicoderesearchcodingbenchmark},
and \textsc{Humanity’s Last Exam}~\citep{phan2025humanitysexam}.
We additionally include \textsc{MMLU-Pro} to assess general academic and
knowledge-intensive reasoning.

Agentic and tool-augmented capabilities are measured using \textsc{TerminalBench},
\textsc{SWE-Bench} (OpenHands) \citep{jimenez2023swe, wang2025openhandsopenplatformai}, 
\textsc{TauBench V2} (airline, retail, telecom) \citep{barres2025tau}, 
and \textsc{BFCL v4} \citep{patil2025bfcl}, each of which provides verifiable reward signals via
unit tests, database state transitions, or structured schema constraints.

Instruction-following and conversational ability are evaluated with \textsc{IFBench},
\textsc{Scale AI Multi-Challenge}, and \textsc{Arena-Hard-V2} \citep{li2024crowdsourced}. 
These benchmarks probe multi-constraint instructions, preference-aligned chat behavior, and faithfulness to user intent. For \textsc{Arena-Hard-V2}, we follow \citet{yang2025qwen3technicalreport} and use GPT-4.1 as judge.

Long-context performance is assessed with \textsc{RULER-100} at 256k, 512k, and 1M tokens \citep{hsieh2024ruler},
together with \textsc{AA-LCR}, evaluating retrieval, stability, and chain-of-thought
coherence over extreme context lengths. RULER-100 is evaluated with reasoning off, whereas AA-LCR is measured with reasoning on.

For multilingual capability, we report results on \textsc{MMLU-ProX} \citep{xuan2025mmlu} and
\textsc{WMT24++} (en$\rightarrow$xx) \citep{deutsch2025wmt24++}, covering a mix of reasoning and translation
settings across multiple high-resource languages.

For comparison with \textsc{GPT-OSS 20B} and \textsc{Qwen3-30B-A3B-Thinking-2507}, we use the officially reported numbers whenever available; if a benchmark is not reported, we take the value from ArtificialAnalysis (AA)\footnote{\url{https://artificialanalysis.ai/}}; and if neither source provides results, we may compute the scores ourselves using the official evaluation protocol.


Table \ref{tab:nano_v3_comparison} presents a comprehensive performance comparison between the three models. \ourmodel shows strong results, surpassing both GPT-OSS 20B and Qwen3-30B-A3B-Thinking-2507 in all categories. On reasoning benchmarks \ourmodel surpasses the Qwen3 model and is competitive with GPT-OSS, which was the previous best model in these categories. In the agentic, chat, and long context categories \ourmodel significantly outperforms both of the other models, demonstrating the strength of our post-training pipeline.
% \todo{eval with tool: is it only python tool?} \\
% \todo{ LCBv6 and LCB range. For now taking published numbers for qwen-3 and gpt-oss} \\
% \todo{Soumye: Model results over stages table on a subset of benchmarks}
% \todo{Key Findings where we put ablations}
% \begin{itemize}
%     \item MoEs beneift from more epochs (in SFT)
%     \item SFT initizlation 5k vs 13k vs 18k
% \end{itemize}

